\section*{一、预习报告}

\subsection*{1. 实验综述}

\textbf{（1）数字示波器的使用。}通过预习，我了解到数字示波器先把待测信号衰减或放大到合适幅度，再由模数转换器（ADC）采样、存储，最后处理并显示波形。为了不失真，采样频率应大于信号最高频率的 2 倍，信号频率一般也不超过示波器带宽的三分之一。触发决定了示波器从哪一点开始采集，只有触发源、触发方式和触发电平设置合适，波形才能稳定显示。

波形稳定后，就可以由屏幕格数读出电压和时间：

\[
V_{p\text{-}p}=D\cdot h_{p\text{-}p},\qquad T=Q\cdot x,\qquad f=\frac1T,
\]

其中 $D$、$Q$ 分别为偏转因数（V/DIV）和时基因数（s/DIV），$h_{p\text{-}p}$、$x$ 为对应的格数。实际测量时，用光标法和自动测量法读取，前者由一对光标的间距读数，后者由 Measure 菜单直接给出。

在此基础上，示波器还能完成三项测量。在 X--Y 模式下，两个频率成简单整数比的正弦信号合成稳定的李萨如图形，满足 $f_y:f_x=N_x:N_y$（$N_x$、$N_y$ 为水平线、竖直线与图形的最多交点数），由已知频率可求出未知频率。在二极管整流电路中，测出输入信号峰峰值 $V_{1p\text{-}p}$ 和输出信号峰值 $V_{2p}$，可得导通电压

\[
V_{\text{导}}=\frac{V_{1p\text{-}p}}{2}-V_{2p}.
\]

在 RC 电路中，输出信号滞后于输入信号，相位差为 $\Delta\varphi=(\Delta t/T)\times360^\circ$，理论值为 $\arctan(2\pi fRC)$。

实验时，先把信号发生器接到 CH1，设置探头比 1×、DC 耦合，调节挡位使波形占据大部分屏幕，再设置触发得到稳定波形。之后依次用光标法和自动法测量峰峰值与频率，在 X--Y 模式下按 $1:1$、$1:2$、$1:3$、$2:1$、$2:3$、$3:4$ 等频率比测量未知频率，并用 $1\,\mathrm{kHz}$、$5\,\mathrm{V}$ 的正弦信号测量二极管导通电压和 RC 电路的相位差。

\textbf{（2）声速的测定。}声速、频率和波长满足 $v=\lambda f$。频率可由信号发生器直接读出，因此关键是测量波长。课件给出温度为 $t$ 时空气中的理论声速

\[
v_{\text{理论}}=331.45\sqrt{1+\frac{t}{273.15\,{}^\circ\mathrm{C}}}\ \mathrm{m/s}.
\]

实验用一对压电陶瓷换能器发射和接收约 $40\,\mathrm{kHz}$ 的超声波，有两种方法测波长。驻波法利用入射波与反射波叠加形成驻波，接收信号幅度每出现一次极大，接收端就移动了 $\lambda/2$；相位比较法把两端信号分别接到示波器的 X、Y 轴，李萨如图形每由一、三象限的直线变为二、四象限的直线，接收端同样移动了 $\lambda/2$。

实验时，先微调信号发生器频率，使接收信号幅度最大，找到谐振频率；再分别用两种方法沿同一方向缓慢移动接收端，各连续记录 12 个位置，同时记下频率和实验前后的室温。数据处理采用逐差法，取 $N=6$：

\[
\lambda_i=\frac{2(x_{N+i}-x_i)}{N},\qquad \bar\lambda=\frac1N\sum_{i=1}^{N}\lambda_i,\qquad v=f\bar\lambda.
\]

不确定度按 $u_v/v=\sqrt{(u_f/f)^2+(u_\lambda/\bar\lambda)^2}$ 合成，再与理论值比较求相对误差。其中 $u_\lambda$ 的 A 类分量取 $\lambda_i$ 平均值的标准偏差；频率和标尺的 B 类分量须以现场仪器的允许误差限为准，本次预习不预设数值。

\subsection*{2. 实验目的}

\begin{enumerate}
\item 了解数字示波器的结构和工作原理，掌握示波器和信号发生器的调节方法。
\item 学会测量信号的幅度、周期和相位差，利用李萨如图形测量未知频率，测量二极管导通电压。
\item 了解声波的特性和干涉现象，用驻波法和相位比较法测定空气中的声速。
\end{enumerate}

\subsection*{3. 实验重点与难点}

\begin{enumerate}
\item 调节触发是获得稳定波形的关键，触发电平须在信号幅度范围内。读数前应让波形尽量占满屏幕但不超出；测相位差时，两路波形要取同类同相点。
\item 观察李萨如图形时，要细调已知信号的频率，等图形基本不再翻转后再读交点数。
\item 测声速前须先调好谐振频率；移动接收端时始终沿同一方向，避免空程差，两换能器也不能靠得过近。测量电路和示波器都需良好接地。
\end{enumerate}
